\section{Model and Scenarios}\label{sec:model}
% Budget 2.5 pages. Table 1.
% generated by experiments/tables/make_table1.py from data/params/model.yaml; do not edit
\begin{table}[t]
\caption{Model parameters as used by the simulator (top) and the two declared modes (bottom). Values are read from the code and listed in \texttt{data/params/model.yaml}.}
\label{tab:params}
\centering\footnotesize\setlength{\tabcolsep}{3pt}
\begin{tabular}{@{}p{1.75cm}p{4.9cm}p{4.9cm}@{}}
\toprule
Agents & \multicolumn{2}{p{9.9cm}}{$m=80$ kg; $v_0\sim\mathcal N(1.34,0.26)$ m/s, $\ge0.5$; $r\sim\mathcal N(0.25,0.03)$ m, $\ge0.15$; $\tau\sim U(0.4,0.6)$ s} \\
Speed--density & \multicolumn{2}{p{9.9cm}}{$v_{\rm eff}=v_0\,\mathrm{clip}[1-e^{-\gamma(1/\rho-1/\rho_{\max})},0.05,1]$, $\gamma=1.913$, $\rho_{\max}=5.4$ ped/m$^2$; $\rho$ = neighbours within $R=3$ m over $\pi R^2$} \\
SFM & \multicolumn{2}{p{9.9cm}}{$A=2000$ N, $B=0.08$ m, $k=1.2\times10^{5}$ kg/s$^2$, $\kappa=2.4\times10^{5}$ kg/(m\,s)} \\
Walls & \multicolumn{2}{p{9.9cm}}{same law per segment, $A_w=A$, $B_w=B$; variants $B_w$, $B_{w,\rm edge}$ (within 0.5 m of the door edge), rounded edge $r=0.1$ m} \\
TTC & \multicolumn{2}{p{9.9cm}}{$k_{\rm ttc}=1.5$, $\tau_0=3$ s, $\tau_{\max}=8$ s} \\
ORCA & \multicolumn{2}{p{9.9cm}}{$\tau_h=5$ s, $\tau_{\rm ORCA}=0.5$ s, speed cap $2v_0$, gate $w_o=1-\sigma(\rho;4,2)$} \\
Crush & \multicolumn{2}{p{9.9cm}}{$k_{\rm crush}=3.6\times10^{5}$, $\kappa_{\rm crush}=4.8\times10^{5}$, gate $w_c=\sigma(\rho;5.5,3)$} \\
Run & \multicolumn{2}{p{9.9cm}}{$\mathrm{d}t=0.01$ s; speed limit $2v_0$; noise $\sigma$ per component, set per experiment; exit at $|x-g|<0.5$ m; deadlock: all $|v|<0.01$ m/s for $10$ s} \\
\midrule
Mode & published-reproduction & corrected-forward \\
\midrule
Integrator & semi-implicit Euler, then clamp; force clip $\pm10^{6}$ N & clamp, then position update; no clip \\
Spawn & uniform, overlaps allowed & rejection, non-overlapping \\
ORCA & as implemented & RVO2-faithful \\
Crush & present in C4 & off \\
Wall corners & sum over segments & convex-union rule \\
\bottomrule
\end{tabular}
\end{table}

%
Agent $i$ is a disc of radius $r_i$ and mass $m_i$ at $\mathbf{x}_i$ with velocity $\mathbf{v}_i$, desired speed $v_{0,i}$ and relaxation time $\tau_i$ (Table~\ref{tab:params} lists all parameters). The equation of motion is $m_i\dot{\mathbf{v}}_i=\mathbf{F}_i$ with
\begin{equation}\label{eq:force}
\mathbf{F}_i=(1-w_{c,i})\,\mathbf{F}^{\rm des}_i+\sum_{j}\mathbf{f}^{\rm sfm}_{ij}+\mathbf{F}^{\rm ttc}_i+w_{o,i}\,\mathbf{F}^{\rm orca}_i+w_{c,i}\,\mathbf{F}^{\rm crush}_i+\sum_{W}\mathbf{f}^{\rm wall}_{iW},
\end{equation}
where $j$ runs over the other active agents within $R=3$\,m and $W$ over the wall segments. The four configurations of the conference paper~\cite{gang2026} switch terms on: C1 is SFM only, C2 adds TTC, C3 adds ORCA, and C4 adds TTC, ORCA and the crush term, which is in the published code but not in the conference description. The code runs in a published-reproduction mode, which reproduces the conference paper's numbers, and a corrected-forward mode, which all new results use (Table~\ref{tab:params}, bottom).

\paragraph{Desired force and density.}
The drive relaxes the velocity towards a speed $v_{{\rm eff},i}$ along the unit vector $\hat{\mathbf e}_i$ from the agent to its goal $\mathbf g$,
\begin{equation}\label{eq:des}
\mathbf{F}^{\rm des}_i=\frac{m_i}{\tau_i}\bigl(v_{{\rm eff},i}\,\hat{\mathbf e}_i-\mathbf v_i\bigr),\qquad
v_{{\rm eff},i}=v_{0,i}\,\mathrm{clip}\!\left[1-e^{-\gamma\,(1/\rho_i-1/\rho_{\max})},\,0.05,\,1\right].
\end{equation}
Here $v_{{\rm eff},i}$ follows Weidmann's speed--density relation~\cite{weidmann1993}, with $\rho_i=N_i/(\pi R^2)$ and $N_i$ the number of other active agents within $R$. In all \TermN\ deadlocks of both modes the blocker, the agent nearest the door, has at most \TermNbMax\ others within $R$ at detection (median \TermNbMedian), so $\rho\le\TermRhoMax$\,ped/m$^2$ and $v_{{\rm eff},i}=v_{0,i}$ to within 1\%. The lateral-bias control is a deterministic tilt of fixed handedness: it replaces $\hat{\mathbf e}_i$ in Eq.~\eqref{eq:des} by $(\hat{\mathbf e}_i+b\,\hat{\mathbf e}^{\perp}_i)/|\hat{\mathbf e}_i+b\,\hat{\mathbf e}^{\perp}_i|$.

\paragraph{Social force and walls.}
For a pair at centre distance $d_{ij}$ with $r_{ij}=r_i+r_j$, unit normal $\mathbf n_{ij}$ pointing from $j$ to $i$, tangent $\mathbf t_{ij}$ and $g(x)=\max(0,x)$,
\begin{equation}\label{eq:sfm}
\mathbf f^{\rm sfm}_{ij}=\bigl[A\,e^{(r_{ij}-d_{ij})/B}+k\,g(r_{ij}-d_{ij})\bigr]\mathbf n_{ij}+\kappa\,g(r_{ij}-d_{ij})\,\bigl[(\mathbf v_j-\mathbf v_i)\cdot\mathbf t_{ij}\bigr]\mathbf t_{ij},
\end{equation}
after~\cite{helbing2000nature}. A wall segment acts through its closest point, at distance $d_{iW}$, with the friction taken against the agent's own velocity:
\begin{equation}\label{eq:wall}
\mathbf f^{\rm wall}_{iW}=\bigl[A_w\,e^{(r_i-d_{iW})/B_w}+k\,g(r_i-d_{iW})\bigr]\mathbf n_{iW}-\kappa\,g(r_i-d_{iW})\,(\mathbf v_i\cdot\mathbf t_{iW})\,\mathbf t_{iW}.
\end{equation}
The published-reproduction mode sums over segments, so a shared end point can act twice; in the corrected-forward mode only the nearer closest point acts when either is the shared end. The door edges are free segment ends and so are sharp.

\paragraph{Time to collision.}
With relative position $\Delta\mathbf x$ and velocity $\Delta\mathbf v$, the contact time $\tau_{ij}$ solves $|\Delta\mathbf x+\Delta\mathbf v\,\tau|=r_{ij}$ and
\begin{equation}\label{eq:ttc}
\mathbf F^{\rm ttc}_i=\sum_j k_{\rm ttc}\,\frac{e^{-\tau_{ij}/\tau_0}}{\tau_{ij}^2}\,\hat{\mathbf n}^{\rm avoid}_{ij},
\end{equation}
where $\hat{\mathbf n}^{\rm avoid}_{ij}$ separates the predicted positions at contact. Pairs count only for $0<\tau_{ij}\le\tau_{\max}$ and $|\Delta\mathbf v|\ge10^{-4}$\,m/s, so $\mathbf F^{\rm ttc}$ is exactly zero when all relative speeds are below $10^{-4}$\,m/s.

\paragraph{ORCA as a force.}
ORCA builds one half-plane of permitted velocities per neighbour, with horizon $\tau_h$, and returns the velocity $\mathbf v^{\rm orca}_i$ nearest to $v_{0,i}\hat{\mathbf e}_i$, with $v_0$, not $v_{\rm eff}$. The model turns it into a force gated off at high density:
\begin{equation}\label{eq:orca}
\mathbf F^{\rm orca}_i=\frac{m_i}{\tau_{\rm orca}}\bigl(\mathbf v^{\rm orca}_i-\mathbf v_i\bigr),\qquad w_{o,i}=1-s(\rho_i;4,2),
\end{equation}
with $s(x;x_0,c)=1/(1+e^{-c(x-x_0)})$. The published solver and a port of the RVO2 library~\cite{vandenberg2011rvo2} differ as listed in Table~\ref{tab:params}. At rest with no binding constraint $\mathbf v^{\rm orca}_i=v_{0,i}\hat{\mathbf e}_i$, so ORCA adds a constant push $w_{o,i}m_iv_{0,i}\hat{\mathbf e}_i/\tau_{\rm orca}$ towards the goal, which we call \emph{ORCA's rest drive}. ORCA's \emph{drive part} is Eq.~\eqref{eq:orca} with $\mathbf v^{\rm orca}_i$ replaced by $v_{0,i}\hat{\mathbf e}_i$; the remainder, zero when no constraint binds, is its \emph{constraint part}. The rest drive belongs to this additive force form; ORCA applied at the velocity level, as in the RVO2 simulator, has none and is outside our scope. The crush term of C4 has the SFM contact form with larger coefficients and a gate $w_{c,i}$ (Table~\ref{tab:params}) that is negligible at deadlock densities.

\paragraph{Integration and noise.}
Since results of social force models depend on the integration scheme~\cite{koster2013}, the corrected-forward mode uses $\mathbf v'=\Pi(\mathbf v+\mathbf F\Delta t/m)$, $\mathbf v'\leftarrow\Pi(\mathbf v'+\boldsymbol\xi)$, $\mathbf x'=\mathbf x+\mathbf v'\Delta t$, where $\Pi$ caps the speed at $2v_0$. The noise $\boldsymbol\xi$ has independent Gaussian components of standard deviation $\sigma$, either fixed per step or scaled as $\sigma\sqrt{\Delta t/0.01\,{\rm s}}$. The published update differs as listed in Table~\ref{tab:params}.

\paragraph{Scenarios.}
In the room bottleneck (Table~\ref{tab:params}) a run is deadlocked when every remaining agent has $|\mathbf v|<0.01$\,m/s for 10\,s without an exit, and the state is then saved as a snapshot. A quench restarts a snapshot under one intervention for at most 120\,s; every intervention is run on the same snapshots, so interventions are compared snapshot by snapshot. Deterministic interventions stop once all speeds stay below $10^{-3}$\,m/s for 10\,s. Two further geometries are used: a bottleneck rebuilt from the CroMa experiments~\cite{fzj_ped20219}, with agent number and corridor length matched run by run (Sect.~\ref{sec:validity}), and an open geometry with counterflow and crossing scenarios \pending{E8, E9 initial-density sweeps; Table~4}.

\paragraph{Statistics.} All bracketed intervals are 95\%. Proportions carry Wilson intervals~\cite{wilson1927} over runs, paired outcomes the exact McNemar test~\cite{mcnemar1947}, times to freeing Kaplan--Meier estimates~\cite{kaplanmeier1958}, and the noise level for half completion, $\sigma_{50}$, a profile-likelihood interval from a logistic fit. Jaccard indices, mean log rate ratios, paired differences and flow ratios carry percentile bootstrap intervals over snapshots, seeds or runs, and CroMa means Student-$t$ intervals. A replay displaces every coordinate by a Gaussian of standard deviation $10^{-9}$\,m; five replays per snapshot and intervention are compared with the unperturbed run on whether and when the blocker is freed. The probes are read jointly as a signature, not as a family of tests, so no correction for multiplicity is applied.
